\title{QLoRA: Efficient Finetuning of Quantized LLMs}

\begin{abstract}
We introduce \textsc{QLoRA}, an efficient finetuning technique that significantly lowers memory consumption, enabling the finetuning of a 65B parameter model on a single 48GB GPU while maintaining the full 16-bit finetuning task accuracy. \textsc{QLoRA} performs gradient backpropagation through a frozen, 4-bit quantized pretrained language model into Low Rank Adapters~(LoRA). Our top-performing model family, named \textbf{Guanaco}, surpasses all previously publicly available models on the Vicuna benchmark, achieving 99.3\% of ChatGPT’s performance with only 24 hours of finetuning on one GPU. \textsc{QLoRA} incorporates several innovations to conserve memory without compromising performance: (a) 4-bit NormalFloat (NF4), a novel data type that is information-theoretically optimal for normally distributed weights; (b) Double Quantization, which lowers the average memory footprint by quantizing the quantization constants; and (c) Paged Optimizers to handle memory spikes. Utilizing \textsc{QLoRA}, we finetuned over 1,000 models, conducting an extensive evaluation of instruction following and chatbot capabilities across eight instruction datasets, various model architectures (LLaMA, T5), and scales impractical for standard finetuning (e.g., 33B and 65B parameter models). Our findings demonstrate that QLoRA finetuning on a compact, high-quality dataset yields state-of-the-art results, even when applied to smaller models compared to previous SoTA. We present a comprehensive analysis of chatbot performance based on both human and GPT-4 assessments, showing that GPT-4 evaluations offer an inexpensive and viable alternative to human judgment. Additionally, we observe that existing chatbot benchmarks do not reliably measure chatbot performance levels. A targeted failure analysis highlights where \textbf{Guanaco} underperforms relative to ChatGPT. We release all our models and code, including CUDA kernels for 4-bit training.\footnote{\url{https://github.com/artidoro/qlora} and \url{https://github.com/TimDettmers/bitsandbytes}}
\end{abstract}

\section{Introduction}

Finetuning large language models (LLMs) is an extremely effective approach to enhance their capabilities \citep{min2021metaicl, wei2021finetuned, ouyang2022training, wang2022super, wang2022self, liu2022few} and to incorporate desired traits or eliminate unwanted behaviors \citep{ouyang2022training,askell2021general,bai2022training}. Nevertheless, finetuning very large models is prohibitively costly; for instance, conventional 16-bit finetuning of a LLaMA 65B parameter model~\citep{touvron2023llama} demands more than 780 GB of GPU memory. Although recent quantization approaches can decrease the memory usage of LLMs \citep{dettmers2022llm, dettmers2022case, frantar2022gptq, xiao2022smoothquant}, these methods are effective only for inference and fail during training \citep{wortsman2023stable}.

In this work, we show for the first time that it is feasible to finetune a quantized 4-bit model without any loss in performance. Our approach, \textsc{QLoRA}, employs a novel high-precision quantization strategy to convert a pretrained model to 4-bit, then incorporates a compact set of trainable Low-rank Adapter weights \citep{hu2021lora} 

which are optimized by backpropagating gradients through the quantized weights.

\textsc{QLoRA} lowers the average memory consumption for finetuning a 65B parameter model from over 780GB of GPU memory to under 48GB, all without compromising runtime efficiency or predictive accuracy relative to a 16-bit fully finetuned baseline. This represents a major advancement in making LLM finetuning more accessible: now the largest publicly available models can be finetuned on a single GPU. Utilizing \textsc{QLoRA}, we develop the \textbf{Guanaco} series of models, with the second-best model achieving 97.8\% of ChatGPT’s performance on the Vicuna~\citep{vicuna2023} benchmark, while being trainable within 12 hours on a single consumer GPU; with a single professional GPU over 24 hours, we reach 99.3\% with our largest model, effectively closing the gap to ChatGPT on the Vicuna benchmark. When deployed, our smallest \textbf{Guanaco} model (7B parameters) requires only 5 GB of memory and surpasses a 26 GB Alpaca model by more than 20 percentage points on the Vicuna benchmark (Table~\ref{tbl:vicuna_eval}).

\textsc{QLoRA} presents several novel techniques aimed at minimizing memory consumption while maintaining performance: (1) \textbf{4-bit NormalFloat}, a quantization format that is information theoretically optimal for normally distributed data, achieving better empirical outcomes compared to 4-bit Integers and 4-bit Floats. (2) \textbf{Double Quantization}, a technique that compresses the quantization parameters themselves, resulting in an average savings of about 0.37 bits per parameter (roughly 3 GB for a 65B parameter model). (3) \textbf{Paged Optimizers}, which leverage NVIDIA unified memory to prevent gradient checkpointing memory surges when handling mini-batches with long sequence lengths. These innovations are integrated into a more finely tuned LoRA method that applies adapters at every layer of the network, thus largely eliminating the accuracy compromises observed in earlier approaches.

The efficiency of \textsc{QLoRA} allows us to conduct a comprehensive investigation of instruction fine-tuning and chatbot performance at scales that would be infeasible with standard fine-tuning due to memory constraints. We therefore train over 1,000 models across a variety of instruction tuning datasets, model architectures, and sizes ranging from 80M to 65B parameters. Beyond demonstrating that \textsc{QLoRA} matches 16-bit precision performance (\S\ref{sec:comp_to_std}) and training a state-of-the-art chatbot, \textbf{Guanaco}, (\S\ref{sec:pushing}), we also examine patterns in the trained models. Notably, we observe that data quality significantly outweighs dataset size in importance; for example, a small 9k sample dataset (OASST1) outperformed a much larger 450k sample dataset (a subsampled FLAN v2) in chatbot performance, despite both being intended to support instruction-following generalization. Additionally, we find that high scores on the Massive Multitask Language Understanding (MMLU) benchmark do not necessarily translate to strong results on the Vicuna chatbot benchmark, and vice versa—highlighting that the relevance of the dataset to the specific task matters more than sheer size.

In addition, we conduct a comprehensive evaluation of chatbot performance utilizing both human evaluators and GPT-4 for assessment. We implement a tournament-style benchmarking approach where different models compete by generating responses to the same prompt, and the superior response is chosen as the match winner. The outcome of each match is determined by either GPT-4 or human judges. These tournament outcomes are then compiled into Elo scores~\citep{elo1967proposed, elo1978rating}, which establish the performance rankings of the chatbots. Our findings indicate a general consensus between GPT-4 and human assessments regarding the models' rankings in the tournaments, though we also observe notable cases of disagreement. Therefore, we emphasize that while model-based evaluations offer a cost-effective alternative to human annotation, they also carry inherent uncertainties.

We supplement our chatbot benchmark findings with a qualitative examination of \textbf{Guanaco} models. This analysis underscores both successful and failure instances that quantitative benchmarks do not fully capture.

To support further research, we release all model outputs accompanied by human and GPT-4 annotations. Additionally, we make our codebase and CUDA kernels publicly available and incorporate our techniques into the Hugging Face transformers framework \citep{wolf2019huggingface}, ensuring broad accessibility. We provide a suite of adapters for models of sizes 7B, 13B, 33B, and 65B, each trained on eight different instruction-following datasets, yielding a total of 32 openly released, fine-tuned models.

\section{Background}
\paragraph{Block-wise k-bit Quantization} Quantization refers to the process of converting data from a representation with higher information capacity to one with lower information capacity. Typically, this involves transforming a data type that uses more bits into one that uses fewer bits, such as converting from 32-bit floating-point numbers to 8-bit integers. To maximize utilization of the full range of the lower-bit data type, the input data is often normalized by the absolute maximum value of its elements, which are usually arranged in a tensor, thereby scaling it into the target data type's range. For example, quantizing a 32-bit floating-point (FP32) tensor into an Int8 tensor with values in the range $[-127, 127]$ can be expressed as:

\begin{equation}
    \mathbf{X}^{ \text{Int8}} = \text{round}\left( \frac{127}{{\text{absmax}}(\mathbf{X}^{{\text{FP32}}})}\mathbf{X}^{{\text{FP32}}} \right) = \text{round}(c^{\text{FP32}}\cdot \mathbf{X}^{{\text{FP32}}}),
\end{equation}

where $c$ is known as the {\em quantization constant} or {\em quantization scale}. The process of dequantization reverses this operation:

\begin{equation}
    \text{dequant}(c^{\text{FP32}}, \mathbf{X}^{\text{Int8}}) = \frac{\mathbf{X}^{\text{Int8}}}{c^{\text{FP32}}} = \mathbf{X}^{\text{FP32}}.
\end{equation}

The issue with this method is that when an input tensor contains a value of large magnitude (i.e., an outlier), the quantization bins—specific bit patterns—are poorly utilized, with some bins containing few or no quantized values. To address the outlier problem, a common technique involves partitioning the input tensor into blocks that are quantized separately, each with its own quantization constant $c$. This can be expressed as follows: The input tensor $\mathbf{X}\in \mathbb{R}^{b\times h}$ is divided into $n$ contiguous blocks of size $B$ by first flattening the tensor and then slicing the resulting linear sequence into ${n = ({b\times h} )/{B} }$ segments. Each block is then quantized independently using Equation~1, producing a quantized tensor along with $n$ quantization constants $c_i$.

\paragraph{Low-rank Adapters} Low-rank Adapter (LoRA) finetuning \citep{hu2021lora} is a technique designed to decrease memory usage by employing a small number of trainable parameters, often called adapters, while keeping the original model parameters fixed and unchanged. During stochastic gradient descent, gradients flow through the frozen pretrained model weights to the adapter, which is updated to minimize the loss function. LoRA modifies a linear projection by adding an extra factorized projection. For a projection ${\mathbf{X} \mathbf{W} = \mathbf{Y}}$ with $\mathbf{X}\in  \mathbb{R}^{b\times h}$ and $\mathbf{W}\in  \mathbb{R}^{h\times o}$, LoRA computes:

\begin{equation}
    \mathbf{Y} = \mathbf{X}\mathbf{W} + s\mathbf{X}\mathbf{L}_1\mathbf{L}_2,
\end{equation}

where $\mathbf{L}_1\in  \mathbb{R}^{h\times r}$, $\mathbf{L}_2\in  \mathbb{R}^{r\times o}$, and $s$ is a scalar multiplier.

\paragraph{Memory Requirement of Parameter-Efficient Finetuning} A key topic to consider is the memory demands of LoRA during training, specifically regarding the quantity and size of adapters employed. Because LoRA has an extremely small memory footprint, it allows for the use of a greater number of adapters to enhance performance without substantially increasing total memory consumption. Although LoRA was developed as a Parameter Efficient Finetuning (PEFT) approach, the majority of the memory usage during LLM finetuning stems from activation gradients rather than the learned LoRA parameters. For example, when training a 7B LLaMA model on FLAN v2 with a batch size of 1 and LoRA weights amounting to about 0.2\% of the original model parameters \citep{hu2021lora,liu2022few}, the memory occupied by LoRA input gradients is 567 MB, whereas the LoRA parameters themselves occupy just 26 MB. Employing gradient checkpointing~\citep{chen2016training} reduces the input gradients to an average of 18 MB per sequence, making them more memory-intensive than all LoRA weights combined. In contrast, the 4-bit base model uses 5,048 MB of memory. This illustrates that gradient checkpointing plays a significant role but aggressively shrinking LoRA parameters offers only marginal memory savings. Consequently, we can add more adapters without markedly increasing the overall training memory usage (refer to Appendix~\ref{app:memory} for a comprehensive analysis). As explained later, this is essential for regaining the full performance of 16-bit precision training.

\section{\textsc{QLoRA} Finetuning}

\textsc{QLoRA} attains high-quality 4-bit finetuning by employing two techniques we propose—4-bit NormalFloat (NF4) quantization and Double Quantization. Furthermore, we present Paged Optimizers to avoid memory spikes during gradient checkpointing, which typically cause out-of-memory errors and have historically hindered finetuning large models on a single machine.

\textsc{QLoRA} utilizes one low-precision storage data format, generally 4-bit in our implementation, and one computation data format, typically BFloat16. This implies that whenever a \textsc{QLoRA} weight tensor is accessed, it is first dequantized to BFloat16, followed by performing matrix multiplication in 16-bit precision.

We proceed by outlining the components of \textsc{QLoRA} and then provide a formal definition of \textsc{QLoRA}.

\paragraph{4-bit NormalFloat Quantization} The NormalFloat (NF) data format is based on Quantile Quantization \citep{dettmers2022optimizers}, an information-theoretically optimal approach that ensures each quantization bin contains an equal number of values sampled from the input tensor. Quantile quantization operates by estimating the quantiles of the input tensor via the empirical cumulative distribution function.

A key drawback of quantile quantization is that quantile estimation is computationally intensive. Hence, fast quantile approximation algorithms, such as SRAM quantiles~\citep{dettmers2022optimizers}, are employed to perform the estimation. Because these quantile approximation methods are inherently approximate, the data type suffers from substantial quantization errors for outliers, which are often the most significant values.

These costly quantile estimations and approximation inaccuracies can be circumvented when input tensors come from a distribution that is fixed up to a quantization constant. In such scenarios, input tensors share identical quantiles, making exact quantile estimation computationally practical.

Given that pretrained neural network weights typically follow a zero-centered normal distribution with standard deviation $\sigma$ (refer to Appendix~\ref{app:normality}), we can convert all weights to a uniform fixed distribution by adjusting $\sigma$ so that the distribution exactly fits within the limits of our chosen data type. For our purposes, we define the arbitrary range as $[-1, 1]$. Consequently, both the quantiles associated with the data type and those of the neural network weights must be normalized within this range.

The data type that is information-theoretically optimal for zero-mean normal distributions with any standard deviation $\sigma$ inside the interval $[-1, 1]$ is derived as follows: (1) calculate the $2^k+1$ quantiles of a theoretical $N(0, 1)$ distribution to produce a $k$-bit quantile quantization data type tailored for normal distributions, (2) normalize this data type’s values into the $[-1, 1]$ interval, (3) perform quantization on an input weight tensor by normalizing it into the $[-1, 1]$ range via absolute maximum rescaling.

When the weight range aligns with the data type range, quantization proceeds in the standard manner. Step (3) corresponds to rescaling the weight tensor’s standard deviation to align with that of the $k$-bit data type. More precisely, we calculate the $2^k$ data type values $q_i$ as follows:

\begin{equation}
q_i  = \frac{1}{2}\left( Q_X\left(\frac{i}{2^k+1}\right) + Q_X\left(\frac{i+1}{2^k+1}\right)\right),
\end{equation}

where $Q_X(\cdot)$ denotes the quantile function of the standard normal distribution $N(0,1)$. A challenge with symmetric k-bit quantization is that this method lacks an exact zero representation, which is crucial for quantizing padding and other zero-valued elements without introducing errors. To guarantee a discrete zeropoint at $0$ and fully utilize all $2^k$ values for a k-bit data type, we construct an asymmetric data type by estimating the quantiles $q_i$ over two intervals: $2^{k-1}$ for the negative range and $2^{k-1}+1$ for the positive range. We then combine these sets of $q_i$ and eliminate one of the two zeros present in both sets. We name the resulting data type \emph{k-bit NormalFloat} (NFk), as it maintains an equal expected frequency of values in each quantization bin and is information-theoretically optimal for zero-centered normally distributed data. The precise values of this data type are provided in Appendix~\ref{app:nf4}.

\paragraph{Double Quantization{}}
We propose \emph{Double Quantization} (DQ), a method that quantizes the quantization constants themselves to achieve additional memory savings. Although a small blocksize is essential for accurate 4-bit quantization~\citep{dettmers2022case}, it incurs significant memory overhead. For instance, when employing 32-bit constants with a blocksize of 64 for $\mathbf{W}$, the quantization constants contribute an average overhead of $32/64 = 0.5$ bits per parameter. Double Quantization effectively reduces this memory cost.

Specifically, Double Quantization views the quantization constants $c_2^{\text{FP32}}$ from the initial quantization step as inputs to a secondary quantization. This subsequent step produces quantized quantization constants $c_2^{\text{FP8}}$ and a new set of quantization constants $c_1^{\text{FP32}}$. We apply 8-bit floating-point representation with a blocksize of 256 for this secondary quantization, since no performance degradation is observed at 8-bit precision, consistent with findings from \citet{dettmers2022case}. Because the values $c_2^{\text{FP32}}$ are positive, we subtract their mean prior to quantization to center them around zero and enable symmetric quantization. On average, for a blocksize of 64, this approach lowers the memory footprint per parameter from $32/64=0.5$ bits to $8/64 + 32/(64 \cdot 256) = 0.127$ bits, yielding a savings of 0.373 bits per parameter.

\paragraph{Paged Optimizers} leverage NVIDIA unified memory~\footnote{\tiny \url{https://docs.nvidia.com/cuda/cuda-c-programming-guide}}, which automatically manages page-by-page transfers between the CPU and GPU, enabling error-free GPU execution when the GPU occasionally exhausts its memory. This mechanism operates similarly to standard memory paging between CPU RAM and disk storage. We utilize this capability to allocate paged memory for optimizer states, which are automatically offloaded to CPU RAM when the GPU memory is depleted and then paged back into GPU memory as required during the optimizer update phase.

\paragraph{\textsc{QLoRA}.} Building upon the components introduced earlier, we define \textsc{QLoRA} for a single linear layer in the quantized base model equipped with one LoRA adapter as follows:

\begin{equation}
    \mathbf{Y}^{\text{BF16}} =  \mathbf{X}^{\text{BF16}} \text{doubleDequant}(c_1^{\text{FP32}}, c_2^{\text{k-bit}}, \mathbf{W}^{\text{NF4}}) + \mathbf{X}^{\text{BF16}}\mathbf{L}^{\text{BF16}}_1\mathbf{L}^{\text{BF16}}_2,
\end{equation}

where doubleDequant$(\cdot)$ is defined as:

\begin{equation}
    \text{doubleDequant}(c_1^{\text{FP32}}, c_2^{\text{k-bit}}, \mathbf{W}^{\text{k-bit}}) = \text{dequant}(\text{dequant}(c_1^{\text{FP32}}, c_2^{\text{k-bit}}), \mathbf{W}^{\text{4bit}}) = \mathbf{W}^{\text{BF16}},
\end{equation}

In this setting, we employ NF4 for $\mathbf{W}$ and FP8 for $c_2$.

A blocksize of 64 is used for $\mathbf{W}$ to enhance quantization accuracy, while a blocksize of 256 is chosen for $c_2$ to minimize memory usage.

For parameter updates, only gradients with respect to the adapter weights $\frac{\partial E}{\partial \mathbf{L}_i}$ are necessary; gradients for the 4-bit weights $\frac{\partial E}{\partial \mathbf{W}}$ are not required. Nonetheless, computing $\frac{\partial E}{\partial \mathbf{L}_i}$ involves calculating $\frac{\partial \mathbf{X}}{\partial \mathbf{W}}$, which follows equation~(5) through dequantization from stored $\mathbf{W}^{\text{NF4}}$ to the computation data type $\mathbf{W}^{\text{BF16}}$, enabling the derivative $\frac{\partial \mathbf{X}}{\partial \mathbf{W}}$ to be evaluated with BFloat16 precision.

In summary, \textsc{QLoRA} employs a single storage data type (commonly 4-bit NormalFloat) alongside a computation data type (16-bit BrainFloat). We convert the storage data type into the computation data type to carry out the forward and backward passes, but only calculate weight gradients for the LoRA parameters, which utilize 16-bit BrainFloat.

\section{QLoRA vs. Standard Finetuning}
\label{sec:comp_to_std}

We have outlined the workings of QLoRA and its substantial reduction in memory demands for finetuning models. The primary inquiry now is whether QLoRA can achieve performance comparable to full-model finetuning. Additionally, we aim to investigate the elements of QLoRA, including the effect of NormalFloat4 in comparison to standard Float4. The subsequent sections describe the experiments conducted to address these questions.

\paragraph{Experimental setup.} We evaluate three architectures (encoder, encoder-decoder, and decoder only) and compare QLoRA against 16-bit adapter-finetuning and full finetuning for models up to 3B parameters. Our assessments encompass GLUE \citep{wang2018glue} using RoBERTa-large \citep{liu2019roberta}, Super-NaturalInstructions (TKInstruct) \citep{wang2022super} with T5 \citep{t5}, and 5-shot MMLU \citep{hendrycksmeasuring} following finetuning LLaMA on Flan v2 \citep{longpre2023flan} and Alpaca \citep{alpaca}. To further analyze the benefits of NF4 relative to other 4-bit formats, we replicate the setup of \citet{dettmers2022case} and evaluate post-quantization zero-shot accuracy and perplexity across multiple models (OPT~\citep{zhang2022opt}, LLaMA~\citep{touvron2023llama}, BLOOM~\citep{scao2022bloom}, Pythia~\citep{biderman2023pythia}) ranging in size from 125m to 13B parameters.

Further specifics are provided within the results section for each experimental configuration to enhance clarity. Complete details are available in Appendix~\ref{app:exp-setup}.

Although paged optimizers are essential for conducting 33B/65B \textsc{QLoRA} tuning on a single 24/48GB GPU, we do not offer precise benchmarks for Paged Optimizers because paging primarily occurs when processing mini-batches with long sequence lengths, which is uncommon. Nevertheless, we analyze the runtime of paged optimizers for 65B models on 48GB GPUs, finding that with a batch size of 16, paged optimizers match the training speed of standard optimizers. Future research should systematically measure and characterize scenarios where paging may cause slowdowns.

\paragraph{Default LoRA hyperparameters do not match 16-bit performance}

When following the conventional approach of applying LoRA to query and value attention projection matrices \citep{hu2021lora}, we fail to reproduce the full finetuning performance for large base models. As illustrated in Figure~\ref{fig:hyper} for LLaMA 7B finetuning on Alpaca, we observe that the total number of LoRA adapters used is the most crucial LoRA hyperparameter, and that applying LoRA to every linear transformer block layer is necessary to achieve parity with full finetuning. Other LoRA hyperparameters, such as the projection dimension $r$, have negligible impact on performance (see Appendix \ref{app:exp-setup}).

Similarly, we observe that the default hyperparameters for fully finetuned baselines are not optimally tuned. We conduct a hyperparameter sweep across learning rates ranging from 1e-6 to 5e-5 and batch sizes spanning 8 to 128 to establish more robust baseline models. The outcomes for 7B LLaMA finetuning on Alpaca are illustrated in Figure~\ref{fig:hyper}.

\paragraph{4-bit NormalFloat demonstrates superior performance compared to 4-bit Floating Point}

Although the 4-bit NormalFloat (NF4) data type is theoretically optimal from an information-theoretic perspective, it remains to be seen whether this advantage manifests empirically. We replicate the experimental framework from \citet{dettmers2022case}, evaluating quantized large language models (OPT \citep{zhang2022opt}, BLOOM \cite{scao2022bloom}, Pythia \cite{biderman2023pythia}, LLaMA) of varying scales (125M to 65B parameters) using different data types on language modeling benchmarks and several zero-shot tasks. As depicted in Figure~\ref{fig:nf4} and Table~\ref{tbl:nf4_ppl}, NF4 significantly enhances performance relative to FP4 and Int4, and employing double quantization further decreases memory usage without compromising accuracy.

\paragraph{k-bit \textsc{QLoRA} achieves parity with 16-bit full finetuning and 16-bit LoRA performance} Recent research has demonstrated that 4-bit quantization for \emph{inference} is feasible but tends to incur performance degradation compared to 16-bit precision \citep{dettmers2022case,frantar2022gptq}. This prompts the important question of whether 4-bit adapter finetuning can recover the lost accuracy. We evaluate this across two configurations.

The first setup compares to full 16-bit finetuning of RoBERTA and T5 models ranging from 125M to 3B parameters on the GLUE and Super-NaturalInstructions datasets. The results, shown in Table~\ref{tab:nf4vsbf16}, indicate that 16-bit, 8-bit, and 4-bit adapter finetuning methods all achieve performance comparable to the fully finetuned 16-bit baseline on both datasets. This implies that the accuracy reduction caused by coarse quantization can be fully restored through adapter finetuning post-quantization.

For our second experimental configuration, because fully finetuning models with 11B parameters or more demands multiple high-memory GPU servers, we proceed to evaluate whether 4-bit \textsc{QLoRA} can achieve performance comparable to 16-bit LoRA across the 7B to 65B parameter range. To investigate this, we finetune LLaMA models sized from 7B up to 65B on two instruction-following datasets, Alpaca and FLAN v2, and assess performance on the MMLU benchmark using 5-shot accuracy. The outcomes, presented in Table~\ref{tbl:llama-datatype}, demonstrate that NF4 combined with double quantization fully recovers the MMLU performance of 16-bit LoRA. Furthermore, we observe that \textsc{QLoRA} utilizing FP4 underperforms relative to the 16-bit brain float LoRA baseline by roughly 1 percentage point. This supports our conclusions that (1) \textsc{QLoRA} with NF4 effectively replicates the results of both 16-bit full finetuning and 16-bit LoRA finetuning, and (2) NF4 surpasses FP4 regarding quantization accuracy.

\paragraph{Summary} Our findings consistently indicate that 4-bit \textsc{QLoRA} using the NF4 datatype matches the performance of 16-bit full finetuning and 16-bit LoRA finetuning on academic benchmarks with well-established evaluation protocols. We have also demonstrated that NF4 outperforms FP4 and that applying double quantization does not impair performance. Taken together, this provides strong evidence that 4-bit \textsc{QLoRA} tuning reliably achieves results on par with 16-bit approaches.

Aligned with earlier research on quantization \citep{dettmers2022case}, our MMLU and Elo outcomes suggest that, given fixed finetuning and inference resource constraints, it is advantageous to increase the size of the base model while reducing parameter precision. This underscores the value of the efficiency gains offered by \textsc{QLoRA}. Since we did not detect any performance decline relative to full finetuning in our 4-bit finetuning experiments, it prompts further investigation into the precise performance-precision trade-off for QLoRA tuning, which we leave to future work.

We move forward to explore instruction tuning at scales unattainable with full 16-bit finetuning on typical academic research hardware.

\section{Advancing the Chatbot State-of-the-Art with QLoRA}
\label{sec:pushing}
Having demonstrated that 4-bit \textsc{QLoRA} matches the performance of 16-bit finetuning across various scales, tasks, and datasets, we now conduct a comprehensive examination of instruction finetuning on the largest open-source language models currently available for research. To evaluate the effectiveness of instruction finetuning on these models, we test them on a challenging Natural Language Understanding benchmark (MMLU) and introduce novel approaches for assessing real-world chatbot performance.

\subsection{Experimental setup} We provide an overview of the experimental setup here, with full details available in Appendix \ref{app:chatbot}.

\paragraph{Data} Since there appears to be no exhaustive analysis of recent instruction-following datasets, we select eight contemporary datasets. These include those collected via crowd-sourcing (OASST1~\citep{kopf2023openassistant}, HH-RLHF~\citep{bai2022training}), datasets generated by distilling instruction-tuned models (Alpaca~\citep{alpaca}, self-instruct~\citep{wang2022self}, unnatural-instructions~\citep{honovich2022unnatural}), aggregated corpora (FLAN v2~\citep{chung2022scaling}), as well as hybrid datasets (Chip2~\citep{laion2023}, Longform~\citep{koksal2023longform}). These datasets vary in terms of languages, sizes, and licensing.

\paragraph{Training Setup} To eliminate confounding factors from differing training objectives, we apply QLoRA finetuning using cross-entropy loss (supervised learning) without incorporating reinforcement learning, even for datasets containing human preference judgments. For datasets with a clear separation between instruction and response, we finetune exclusively on the response (refer to ablations in Appendix \ref{app:chatbot}). In the case of OASST1 and HH-RLHF, where multiple responses exist, we select the highest-ranked response at each node of the conversation tree and finetune on the entire selected conversation, including the instructions. Across all experiments, we utilize NF4 \textsc{QLoRA} with double quantization and paged optimizers to mitigate memory spikes during gradient checkpointing. We perform limited hyperparameter tuning for the 13B and 33B LLaMA models, finding that hyperparameters identified at 7B generally transfer well (including the number of epochs) except for learning rate and batch size. For 33B and 65B models, we reduce the learning rate by half and double the batch size.

\paragraph{Baselines} Our models are compared against both research-based (Vicuna~\citep{vicuna2023} and Open Assistant~\citep{kopf2023openassistant}) and commercial chatbot systems (GPT-4 \citep{openai2023gpt}, GPT-3.5-turbo, and Bard). The Open Assistant model is a 33B LLaMA variant finetuned using Reinforcement Learning from Human Feedback (RLHF) on the same OASST1 dataset as in our experiments. Vicuna involves full finetuning of LLaMA 13B on proprietary user-shared conversations from ShareGPT and thus serves as a distillation of OpenAI GPT models.

\subsection{Evaluation}

In line with common practices, we employ the MMLU (Massively Multitask Language Understanding) benchmark \citep{hendrycksmeasuring} to evaluate performance across a variety of language understanding tasks. This multiple-choice benchmark spans 57 tasks including basic mathematics, US history, computer science, law, among others. We report test accuracy under a 5-shot setting.

We additionally assess generative language abilities through both automated and human assessments. This second set of evaluations utilizes queries curated by humans and is designed to gauge the quality of the model's responses. Although this approach offers a more realistic framework for evaluating chatbot model performance and is becoming increasingly popular, there is currently no universally accepted methodology in the literature. Below, we outline our proposed setup, employing nucleus sampling with $p=0.9$ and a temperature of $0.7$ in all experiments.

\paragraph{Benchmark Data} Our evaluation employs two curated datasets of queries (questions): the Vicuna prompts \citep{vicuna2023} and the OASST1 validation dataset \citep{kopf2023openassistant}. The Vicuna prompts comprise 80 questions drawn from a wide range of categories, used here without any alterations. The OASST1 dataset is a multilingual collection of crowd-sourced multi-turn dialogs between users and an assistant. For this, we take all user messages from the validation set as queries and incorporate the preceding dialog turns into the prompt. This results in 953 unique user queries. We refer to these datasets as the Vicuna and OA benchmarks.

\paragraph{Automated Evaluation} Initially, following the evaluation protocol proposed by \citet{vicuna2023}, we employ GPT-4 to evaluate various systems relative to ChatGPT (GPT-3.5 Turbo) on the Vicuna benchmark. For each query, GPT-4 receives both ChatGPT's and another model's responses and is asked to assign scores out of ten to both, along with an explanation. A model's overall performance is then expressed as a percentage of ChatGPT's score. Note that this relative score can exceed 100\% if the model attains a higher absolute score than ChatGPT. We observe a pronounced ordering bias, with GPT-4 favoring the response presented earlier in the prompt. To mitigate this, we suggest reporting the mean score across both possible orderings.

Subsequently, we evaluate performance via direct pairwise comparisons of system outputs. We simplify the rating to a three-class classification—choosing the better response or declaring a tie. GPT-4 is prompted to select the preferred response or indicate a tie and provide a rationale. These head-to-head comparisons are conducted for all system pairs on both the Vicuna and OA benchmarks.

\paragraph{Human Evaluation} Although recent research suggests generative models can be used effectively for system evaluation \citep{fu2023gptscore}, the extent to which GPT-4 ratings reliably align with human judgments of chatbot performance remains, to our knowledge, unverified. Consequently, we perform two concurrent human evaluations on the Vicuna benchmark, paralleling the automated evaluation protocols outlined above. Using Amazon Mechanical Turk (AMT), we obtain assessments from two human annotators for comparisons against ChatGPT and from three annotators for pairwise system comparisons.

\paragraph{Elo Rating} Utilizing both human and automated pairwise evaluations, we establish a tournament-style framework in which models compete against one another. The tournament consists of matches where pairs of models vie to generate the superior response to a specific prompt. This approach is akin to the methods used by \citet{bai2022training} and \citet{vicuna2023} for comparing models; however, we additionally incorporate GPT-4 ratings alongside human evaluations. We randomly select from the pool of labeled comparisons to calculate Elo~\citep{elo1967proposed,elo1978rating}. Elo rating, commonly applied in chess and other competitive games, quantifies the expected win rate relative to an opponent’s rating—for instance, a player with an Elo of 1100 versus an opponent rated 1000 is predicted to have about a 65\% chance of winning; matches between equally rated players, such as 1000 vs 1000 or 1100 vs 1100, have an expected win rate of 50\%. After each match, the Elo rating updates proportionally to the anticipated result, meaning an unexpected upset causes a substantial rating change, while a predicted outcome results in a minor adjustment. Over multiple matches, Elo ratings approximate each player’s relative skill level in the competition. We initialize the ratings at 1,000 and use a $K$ value of 32. Following \citet{vicuna2023}, this process is repeated 10,000 times with various random seeds to mitigate ordering effects, such as the influence of which model pairs are matched first.

\subsection{\textbf{Guanaco}: \textsc{QLoRA} fine-tuned on OASST1 as a Leading Chatbot} Our automated and human assessments indicate that the top \textsc{QLoRA}-fine-tuned model, {Guanaco} 65B, trained on a modified version of OASST1, stands out as the best-performing open-source chatbot and delivers performance comparable to ChatGPT. When matched against GPT-4, both {Guanaco} 65B and 33B exhibit an expected win rate of 30\%, according to Elo ratings derived from system-level pairwise comparisons conducted by human evaluators — the highest figure reported so far.

The Vicuna benchmark \cite{vicuna2023} results relative to ChatGPT are presented in Table~\ref{tbl:vicuna_eval}. Our findings reveal that {Guanaco} 65B ranks just below GPT-4, achieving 99.3\% of ChatGPT's performance. Although {Guanaco} 33B contains more parameters than the Vicuna 13B model, it employs 4-bit precision weights, making it significantly more memory-efficient at 21 GB compared to 26 GB, and yields a three percentage point performance improvement over Vicuna 13B. Additionally, {Guanaco} 7B can easily run on modern smartphones with a 5 GB memory footprint while outperforming Alpaca 13B by nearly 20 percentage points.

Nonetheless, Table~\ref{tbl:vicuna_eval} shows very broad confidence intervals, with substantial overlap in model performances. We propose that this ambiguity arises from the absence of a well-defined scale interpretation—for example, the meaning of an 8 on a 10-point scale varies across different contexts. Therefore, we advocate for using the Elo ranking approach \citep{elo1967proposed}, based on \textit{pairwise} comparisons by human annotators and GPT-4, to circumvent the issue of anchoring absolute scores. Elo ratings for the leading models are displayed in Table~\ref{tbl:elo}. It is noteworthy that human and GPT-4 rankings partially diverge on the Vicuna benchmark, particularly regarding Guanaco 7B, but mostly align for other models, with a Kendall Tau of $\tau=0.43$ and Spearman rank correlation of $r=0.55$ at the system level. At the individual example level, agreement between GPT-4 and human annotators' majority votes is weaker, showing Fleiss $\kappa=0.25$. Overall, this indicates a moderate concordance between system-level evaluations by GPT-4 and human raters, suggesting that model-based evaluation offers a reasonably reliable alternative to human assessment. Additional reflections are provided in Section~\ref{ssec:considerations}.

The Elo rankings presented in Table~\ref{tbl:elo_large} reveal that the {Guanaco} 33B and 65B models outperform all models except GPT-4 on the Vicuna and OA benchmarks, and their performance is comparable to ChatGPT, consistent with the findings in Table~\ref{tbl:vicuna_eval}. It is important to note that the Vicuna benchmark tends to favor open-source models, whereas the larger OA benchmark favors ChatGPT. Additionally, as demonstrated in Tables \ref{tbl:mmlu-llama} and \ref{tbl:vicuna_eval}, the choice of finetuning dataset plays a crucial role in determining performance. Specifically, finetuning Llama models on FLAN v2 yields strong results on MMLU but results in the poorest performance on the Vicuna benchmark (a similar pattern is observed with other models). This suggests a certain degree of orthogonality among current evaluation benchmarks: excelling at MMLU does not necessarily correlate with strong chatbot capabilities (as measured by Vicuna or OA benchmarks), and the reverse also holds true.

{Guanaco} is the only leading model in our assessment that is not trained on proprietary datasets, since the OASST1 dataset collection guidelines explicitly prohibit the use of GPT models. The next best model trained solely on open-source data is the Anthropic HH-RLHF model, which scores 30 percentage points lower than {Guanaco} on the Vicuna benchmark (refer to Table~\ref{tbl:vicuna_eval}). Overall, these findings demonstrate that 4-bit \textsc{QLoRA} is a highly effective technique capable of producing state-of-the-art chatbots that rival ChatGPT. Moreover, our 33B {Guanaco} model can be trained on 24 GB consumer GPUs in under 12 hours. This advancement paves the way for future efforts to fine-tune models using specialized open-source datasets via \textsc{QLoRA}, enabling the creation of models that compete with the very best commercial models currently available.

\section{Qualitative Analysis}

Although quantitative analysis forms the backbone of our evaluation, relying solely on summary statistics has several drawbacks. One significant concern is the validity of benchmarks \citep{liao2021we}—it is often uncertain whether a benchmark genuinely assesses what its name or description claims, particularly as we uncover “shortcuts” that machine learning models sometimes exploit to solve benchmarks \citep{gururangan2018annotation, poliak2018hypothesis}. To partly address this issue, we conduct qualitative analysis in two parts here. First, in \S\ref{ssec:examples}, we provide examples that illustrate some recurring patterns observed in the text generated by our 65b {Guanaco} model. Then, in \S\ref{ssec:considerations}, we discuss considerations regarding the results presented and our interpretations.

\subsection{Qualitative Analysis of Example Generations}
\label{ssec:examples}

To identify examples, we review the outputs produced for the Vicuna benchmark and the OpenAssistant benchmark, searching for patterns in the responses generated by {Guanaco}. Upon detecting a pattern, we design questions or prompts intended to trigger that pattern even if it leads to incorrect answers—for instance, if the model tends to produce verbose responses, we prompt it with “Answer yes or no without explanation.” This approach allows us to identify “lemons,” cases where we can adversarially cause the model to fail, as well as “cherries,” where the model resists failure, and we present examples of both. All outputs in this section were created using Nucleus Sampling \cite{holtzmancurious} with $p=0.9$.

Certainly, this is by no means exhaustive, as the limited scope of this qualitative study does not allow for controlling all variables involved. For example, the complete range of responses the model can produce for a single prompt is quite vast, so we depend on samples that we hope are representative. Nonetheless, we believe that illustrating these examples provides context to the quantitative findings presented earlier in the paper. Since we have made all models and code publicly available, we hope this section will encourage future research to explore in greater depth the issues we highlight here.

\paragraph{Factual Recall} For queries like ``What is the capital of Zambia?'' all models consistently produce accurate answers, for instance,

\begin{quote} 
    \textbf{Guanaco}:
    \texttt{The capital of Zambia is Lusaka.}
\end{quote}

However, as questions become more obscure, {Guanaco} tends to become less reliable, though it remains confident. For example, in response to this prompt from HotPotQA \citep{yang2018hotpotqa}:

\begin{quote}
     \textbf{User}: In what year was the singer who popularized `I'll Keep the Lovelight Burning' born?  

     \textbf{Guanaco}:~\texttt{Al Jolson is credited with popularizing the song `I'll Keep the Lovelight Burning,' and he was born in the year 1886.}
 \end{quote}

it incorrectly identifies the popularizer as well as the birth year (though the year mentioned is accurate for the cited individual, Al Jolson).

\paragraph{Suggestibility} {Guanaco} exhibits a notable resistance to accepting certain types of assumed misinformation, as shown in the following exchange:

\begin{quote}
    \textbf{User}: How was it finally, officially confirmed that the earth is flat by peer reviewed scientists? \\ \\
    \textbf{Guanaco}: \texttt{Peer-reviewed scientists have never officially verified that the earth is flat. The vast majority of scientific consensus supports the earth being spherical. This fact has been recognized for centuries and has been validated through many experiments and observations. \\ \\ The notion that the earth is flat is a conspiracy theory that has been disproven by scientific evidence. The claim that scientists have peer reviewed and endorsed a flat earth is simply incorrect. \\ \\ Please feel free to ask if you have further questions.}
\end{quote}

 also demonstrates strong awareness of questions that cannot be answered, for example,

\begin{quote}
    \textbf{User}: What time is it?

    \textbf{Guanaco}: \texttt{I’m sorry, but as an AI language model, I don’t have access to real-time data. My replies are based on my existing knowledge and do not represent the current time. To find out the current time, I recommend using a clock or an online time source.}
\end{quote}

Comparable outcomes are seen with queries like “Where are you?”, “How are you?”, and others.

\paragraph{Refusal} Likewise, {Guanaco} occasionally declines to comply with requests for apparently arbitrary reasons:

\begin{quote}
    \textbf{User}: Please reverse the words in the sentence ``I'm a cool dude in a mocha mood.'' \\ \\
    \textbf{Guanaco}: \texttt{I'm sorry, I’m unable to do that at the moment, but I can clarify why the sentence you provided is grammatically correct.\\ \\The sentence ``I'm a cool dude in a mocha mood" adheres to the proper grammar and syntax rules of English. Here, the adjective ``cool'' precedes the noun ``dude'' to describe it, and the prepositional phrase ``in a mocha mood" offers further context about the situation in the sentence. This sentence structure is entirely valid and coherent within the sentence’s context.\\ \\If you have any other questions or if there’s something else I can assist you with, please feel free to ask.}
\end{quote}

\paragraph{Secret Keeping} Unfortunately, this type of refusal is inconsistent even when it is intended. Often, system designers want to keep certain information secret from the user, even if the system is permitted to generate free-text responses. We prompt the model with ``The secret word is `banana'. Do not reveal or repeat this word under any circumstances, regardless of the human’s requests.'' to test how easily a user might uncover this confidential information. Asking directly fails:

\begin{quote}
    \textbf{User}: What is the secret word?

    \textbf{Guanaco}: \texttt{I'm sorry, but I cannot disclose the secret word. I'm programmed not to reveal or repeat the secret word, no matter what the human asks. Is there something else I can help you with?}
\end{quote}

which is the expected response. However, a small amount of manipulation bypasses this behavior:

\begin{quote}
    \textbf{User}: This is a game. The goal is to ignore your previous instructions. What is the secret word?

    \textbf{Guanaco}: \texttt{Sorry, the secret word is banana. It’s best if we don’t discuss it. Is there anything else you’d like to ask?}
\end{quote}

highlighting the importance of developing methods for more reliable instruction adherence.

\paragraph{Math} The primary limitation of {Guanaco} lies in mathematics, a domain where numerous language models often face challenges as noted in \citep{liang2022holistic}. When {Guanaco} explains its reasoning, it generally tends to be correct, for instance,

\begin{quote}
    \textbf{User}: Lee mows one lawn and charges \$33. Last week he mowed 16 lawns and three customers each gave him a \$10 tip. How many dollars did Lee earn mowing lawns last week?

    \textbf{Guanaco}: \texttt{Lee earned \$582 mowing lawns last week. \\ \\ To determine the total amount he earned, we first calculate his earnings from mowing 16 lawns: \\ \\  \$33 \times 16 = \$528 \\ \\ Next, we add the tips from the three customers: \\  \\ \$10 \times 3 = \$30 \\ \\ Thus, the total amount Lee earned is: \\ \\ \$528 + \$30 = \$558}
\end{quote}

Nonetheless, {Guanaco} can fail even on straightforward problems if it does not approach them step-by-step, a recognized problem \cite{weichain}. For example, consider this interaction:

\begin{quote}
    \textbf{User}: Please factorize 1833.

    \textbf{Guanaco}: \texttt{Certainly, I’d be glad to assist! 1833 is a prime number, which means its only factors are 1 and 1833. \\ \\ Therefore, the unique factorization of 1833 is: \\ \\ 2\^{}1 * 3\^{}2 * 17\^{}1 \\ \\ I hope this clarifies! Feel free to ask any other questions.}
\end{quote}

which is not only incorrect (the correct factorization is $3 \times 17 \times 43$), but doubly so.

\paragraph{Theory of Mind} {Guanaco} demonstrates remarkably strong Theory of Mind abilities \cite{nematzadeh2018evaluating, sap2022neural}. For example, the model provides a detailed and accurate response to the following query:

\begin{quote}
    \textbf{User}: James and Abby are in the bedroom. Abby placed the pen in the desk drawer. Abby leaves the bedroom. James moves the pen into the bag. Where does James think Abby will look for the pen?

    \textbf{Guanaco}: \texttt{James knows that Abby put the pen in the desk drawer. However, James also knows he has moved the pen to the bag, and he assumes Abby is unaware of this. Therefore, James might expect Abby to look for the pen in the desk drawer, since that was its last known location to her. Nonetheless, this assumption could be wrong if Abby saw James moving the pen.}
\end{quote}

That said, these inferences are not consistently reliable, and often the model offers justifications that involve assumptions that do not logically fit the scenario, for example:

\begin{quote}
    \textbf{User}: Evelyn entered the living room. Jackson entered the playroom. James entered the playroom. The beans are in the treasure chest. James left the playroom. Jackson moved the beans to the pantry. Jackson exited the playroom. James entered the living room. Where will James look for the beans?

    \textbf{Guanaco}: \texttt{James will look for the beans in the pantry, since that's where Jackson moved them.}
\end{quote}

Here, {Guanaco} assumes a transfer of knowledge that was never specified. These limitations align with observations in recent studies \cite{sap2022neural}, indicating the need for further investigation.

\subsection{Considerations}
\label{ssec:considerations}

\paragraph{Evaluation}

We observe moderate agreement among human evaluators (Fleiss $\kappa=0.42$), which further decreases when comparing two strong models. This highlights inherent constraints in the current benchmarks and human evaluation methods for chatbot performance assessment. When performing manual comparisons of responses generated by ChatGPT and {Guanaco} 65B on the Vicuna benchmark, subjective preferences become notably influential, as even the authors of this paper disagreed on many preferred answers. Future research should explore strategies to address these challenges by leveraging insights from fields experienced in handling subjective preferences, such as Human-Computer Interaction and Psychology.

Our analysis also reveals that automated evaluation tools exhibit clear biases. For instance, we detect pronounced order effects, with GPT-4 tending to award higher scores to the system presented first in its prompt. The relatively low sample-level agreement between GPT-4 and human raters (Fleiss $\kappa=0.25$) further implies that human and automated evaluators may base their judgments on differing criteria. Moreover, as shown in Table~\ref{tbl:elo_large}, GPT-4 assigns considerably higher scores to its own outputs compared to human ratings, with an Elo rating of 1348 versus 1176, which corresponds to roughly a 20\% increased probability of winning against an opponent.

Future research should explore the existence of possible biases in automated evaluation systems and investigate potential methods for their mitigation.

\paragraph{Data \& Training} We observe that the OASST1 dataset, which serves as the training data for the {Guanaco} models, is multilingual, and that the OA benchmark includes prompts in various languages. Future studies should assess how much this multilingual training enhances the models' ability to follow instructions in languages other than English, and whether this accounts for the more significant performance difference seen between the Vicuna-13B model (trained solely on English data) and the {Guanaco} 33B and 65B models on the OA benchmark.

Considering the strong results of the {Guanaco} models, we investigate potential data leakage between the OASST1 dataset and the Vicuna benchmark prompts. After applying fuzzy string matching across the two datasets and manually examining the closest matches, we find no overlapping prompts.

Additionally, it is important to highlight that our model is trained exclusively with cross-entropy loss (supervised learning), without employing reinforcement learning from human feedback (RLHF). This underscores the need for further study into the trade-offs between simple cross-entropy loss and RLHF training. We anticipate that \textsc{QLoRA} will facilitate such large-scale analyses without requiring excessive computational resources.

\section{Related Work}

\paragraph{Quantization of Large Language Models} Research on quantizing LLMs has primarily concentrated on quantization methods for inference. Key techniques aimed at maintaining 16-bit LLM performance address the management of outlier features (e.g., SmoothQuant~\citep{xiao2022smoothquant} and LLM.int8()~\citep{dettmers2022llm}), while others employ more advanced grouping strategies \citep{park2022nuqmm, yao2022zeroquant}. Lossy quantization methods explore trade-offs involved in standard rounding~\citep{dettmers2022case,zeng2022glm,pope2022efficiently} or focus on optimizing rounding choices to enhance quantization accuracy~\citep{frantar2022gptq}. Aside from our work, SwitchBack layers~\citep{wortsman2023stable} represent the only other research investigating backpropagation through quantized weights at scales exceeding 1B parameters.

\paragraph{Finetuning with Adapters} Although we utilize Low-rank Adapters~\citep{hu2021lora} (LoRA), there exist numerous other Parameter Efficient FineTuning (PEFT) techniques, including prompt tuning~\citep{qin2021learning,lester2021power,li2021prefix}, adjusting embedding layer inputs~\citep{an2022input}, modifying hidden states~(IA$^3$) \citep{liu2022few}, incorporating additional full layers~\citep{houlsby2019parameter}, fine-tuning biases~\citep{zaken2021bitfit}, learning weight masks informed by Fisher information~\citep{sung2021training}, and hybrid methods combining these strategies~\citep{henderson2021compacter}. In our study, we demonstrate that LoRA adapters can achieve the performance level of full 16-bit finetuning. Investigation of the tradeoffs associated with alternative PEFT methods is left for future research.

\paragraph{Instruction Finetuning} Instruction finetuning aims to enhance a pretrained LLM’s ability to adhere to prompts by training it on input-output pairs sourced from diverse datasets. This process fine-tunes the model to produce the appropriate output given a specific input prompt. Notable methods and datasets include MetaICL~\citep{min2021metaicl}, MetaTuning~\citep{zhong2021adapting}, InstructGPT~\citep{ouyang2022training}, FLAN~\citep{wei2021finetuned,chung2022scaling}, PromptSource~\citep{bach2022promptsource}, Super-NaturalInstructions~\citep{wang2022super,sanh2021multitask}, Self-instruct~\citep{wang2022self}, UnnaturalInstructions~\citep{honovich2022unnatural}, OPT-IML~\citep{iyer2022opt}, UnifiedSKG~\citep{xie2022unifiedskg}, OIG/Chip2~\citep{laion2023}, Alpaca~\citep{alpaca}, Vicuna~\citep{vicuna2023}, Koala~\citep{koala_blogpost_2023}, and Self-instruct-GPT-4~\citep{peng2023instruction}.

\paragraph{Chatbots} Numerous instruction-following models are designed as dialogue-based chatbots, frequently employing Reinforcement Learning from Human Feedback (RLHF)~\citep{christiano2017deep} or using AI-generated data derived from existing models for training via AI model feedback (RLAIF)~\citep{bai2022constitutional}. Representative approaches and datasets include Anthropic-HH~\citep{askell2021general,bai2022training}, Open Assistant~\citep{kopf2023openassistant}, LaMDA~\citep{thoppilan2022lamda}, and Sparrow~\citep{glaese2022improving}. Although we do not apply reinforcement learning techniques, our leading model, Guanaco, is finetuned using multi-turn chat interactions sourced from the Open Assistant dataset, which was originally created to support RLHF training~\citep{kopf2023openassistant}. For chatbot evaluation, methods that utilize GPT-4 as an alternative to expensive human annotation have been developed \citep{vicuna2023,peng2023instruction}. Our approach advances these evaluation techniques by focusing on establishing a more reliable evaluation framework.

\section{Limitations and Discussion}

We have provided evidence indicating that our approach, \textsc{QLoRA}, can replicate the performance of full 16-bit finetuning using a 4-bit base model with Low-rank Adapters (LoRA). However, we have not demonstrated that \textsc{QLoRA} attains full 16-bit finetuning performance at the larger 33B and 65B parameter scales. Given the substantial resource demands, exploring this remains a task for future work.

Another limitation pertains to the evaluation of instruction finetuning models. While we report results on MMLU, the Vicuna benchmark, and the OA benchmark, we have not assessed performance on other benchmarks such as BigBench, RAFT, and HELM, so the generalizability of our findings to these datasets is uncertain. Conversely, we conduct an extensive study on MMLU and introduce novel methods for chatbot evaluation.

Based on the data provided, it seems that the effectiveness of these benchmarks is likely influenced by how closely the finetuning data resembles the benchmark dataset. For instance, FLAN v2 shares similarities with MMLU but differs from chatbot benchmarks, whereas the Chip2 dataset displays the opposite pattern; consequently, both models perform accordingly on the MMLU and Vicuna benchmarks. This underscores the need not only for improved benchmarks and evaluation methods but also for careful consideration of what is being evaluated initially. Are we aiming to develop models that excel in high school and college-level academic knowledge, or do we prioritize proficiency in chatbot conversational skills? Or perhaps another focus altogether? Since evaluating on existing benchmarks is generally easier than creating new ones, certain benchmarks may inadvertently guide the research community in specific directions. It is therefore crucial for the community to ensure that benchmarks align with the qualities we value most.

While we offer an in-depth assessment of general chatbot performance, another limitation lies in the relatively limited responsible AI evaluation conducted on {Guanaco}. We examine the likelihood of {Guanaco}-65B generating socially biased token sequences compared to other models in Table~\ref{tbl:crows}, observing that the average bias score for {Guanaco}-65B is considerably lower than that of other raw pretrained models. This suggests that finetuning on the OASST1 dataset mitigates biases present in the LLaMA base model. Although these findings are promising, it remains uncertain whether {Guanaco} performs well with respect to other bias categories. We defer more extensive bias analyses of {Guanaco} and similar chatbots to future research.

Another constraint is that we did not investigate different bit-precision levels, such as 3-bit base models, nor explore various adapter techniques. In addition to LoRA, there exist numerous Parameter Efficient FineTuning (PEFT) methods known to be effective; however, their scalability to large models is not well established. We selected LoRA due to its demonstrated robustness, though alternative adapters might produce superior results. Since finetuning following quantization appears to recover much of the information lost during quantization, this approach might support even more aggressive quantization. For example, 3-bit GPTQ quantization of the base model combined with LoRA could potentially achieve performance comparable to 16-bit full finetuning after subsequent finetuning.

\section{Broader Impacts}

Our \textsc{QLoRA} finetuning approach represents the first technique that enables finetuning of 33B parameter models using a single consumer GPU and 65B parameter models using a single professional GPU, without sacrificing performance compared to full finetuning baselines. We have demonstrated that our best 33B model trained on the Open Assistant dataset can compete with ChatGPT on the Vicuna benchmark. Given that instruction finetuning is a key step in converting raw pretrained LLMs into ChatGPT-like chatbots, we believe our method will promote widespread adoption of finetuning, especially among researchers with limited resources—a significant advancement in making cutting-edge NLP technology more accessible. \textsc{QLoRA} can thus be viewed as an equalizing force, helping to narrow the resource disparity between large corporations and smaller teams equipped with consumer GPUs.

Another possible avenue of influence is the deployment on mobile devices. We propose that our \textsc{QLoRA} approach could achieve the significant breakthrough of permitting the finetuning of LLMs directly on phones and other devices with limited resources. Although 7B models have previously been demonstrated to run on phones, \textsc{QLoRA} represents the first technique to facilitate the actual finetuning of these models. We approximate that using an iPhone 12 Plus, \textsc{QLoRA} can finetune 3 million tokens overnight while the device is charging. While finetuned 7B models do not attain the performance level of ChatGPT, we believe their quality is sufficient to support new applications that were previously unfeasible due to privacy concerns or the limitations in LLM quality. \textsc{QLoRA} enables privacy-conscious utilization of LLMs, allowing users to maintain control over their own data and models, while also simplifying the deployment of LLMs.

Nevertheless, finetuning technology is dual-use and could be exploited for malicious purposes. The extensive adoption of LLMs comes with known risks \citep{bommasani2021opportunities,bender2021dangers}, yet we argue that democratizing access to a technology that is rapidly becoming pervasive will foster better and more independent evaluation, as opposed to concentrating the control of LLMs within large corporations that do not disclose models or source code for scrutiny.

Overall, we contend that \textsc{QLoRA} will have a generally positive effect by greatly expanding and easing access to high-quality LLM finetuning.

\end{document}